\section{Methodology}

\subsection{Overview}
Our methodology combines quantitative data analysis with qualitative LLM-based reasoning. The system pipeline consists of three distinct stages: Data Acquisition and Preprocessing, Prompt Engineering and Context Construction, and Risk-Adjusted Decision Logic.

\subsection{Data Acquisition and Preprocessing}
\subsubsection{Data Sources}
We utilize the Yahoo Finance API as our primary data source due to its reliability and comprehensive coverage of global equities. For any given ticker symbol $S$, we fetch two datasets:
\begin{itemize}
    \item $Q_t$: The real-time quote at time $t$, including current price $P_t$, volume $V_t$, and daily change $\Delta P_t$.
    \item $H_{t-30 \dots t}$: A historical time-series of Daily OHLCV (Open, High, Low, Close, Volume) data for the past 30 trading days.
\end{itemize}

\subsubsection{Data Normalization}
To ensure the LLM perceives relative movements rather than absolute values (which can vary by orders of magnitude between stocks), we normalize the historical data. We calculate the daily percentage return $R_i$ for each day $i$:
\begin{equation}
    R_i = \frac{P_i - P_{i-1}}{P_{i-1}} \times 100
\end{equation}
This transformation allows the model to analyze volatility and momentum independent of the nominal share price.

\subsection{Mathematical Model of Risk}
A core component of our "Low Risk" strategy is the quantification of volatility. While the LLM provides a qualitative assessment, we compute quantitative metrics to validate the "safety" of a trade.

\subsubsection{Volatility ($\sigma$)}
We define the 30-day historical volatility $\sigma$ as the standard deviation of the daily log returns:
\begin{equation}
    \sigma = \sqrt{\frac{1}{N-1} \sum_{i=1}^{N} (r_i - \bar{r})^2}
\end{equation}
where $r_i = \ln(P_i / P_{i-1})$ and $N=30$. Stocks with a $\sigma$ exceeding a predefined threshold (e.g., 0.05) are automatically flagged as "High Risk" before being sent to the LLM.

\subsubsection{Sharpe Ratio}
To assess the risk-adjusted return potential, we estimate a rolling Sharpe Ratio:
\begin{equation}
    S_a = \frac{E[R_a - R_f]}{\sigma_a}
\end{equation}
where $R_a$ is the asset return and $R_f$ is the risk-free rate (approximated by the 10-year Treasury yield).

\begin{lstlisting}[language=JavaScript, caption=Prompt Template]
const prompt = `
  Analyze the following stock data for ${symbol}.
  Current Price: ${quote.regularMarketPrice}
  Recent History (last 30 days): ${JSON.stringify(recentHistory)}

  Based on this data, determine if the trend suggests a "BUY" or "SELL" or "HOLD".
  Provide a recommendation and a short textual explanation (max 3 sentences) explaining the trend.
  
  Format the response as JSON:
  {
    "recommendation": "BUY" | "SELL" | "HOLD",
    "explanation": "string"
  }
`;
\end{lstlisting}

\subsection{Risk Assessment Logic}
The "Low Risk" aspect of our strategy is enforced through the "HOLD" recommendation. The model is instructed (via the few-shot examples inherent in its training and the prompt's phrasing) to recommend "HOLD" when the trend is ambiguous or volatility is high. This conservative bias is crucial for retail investors.

\subsection{Prompt Engineering Framework}
The quality of the LLM's output is directly dependent on the quality of the prompt. We employ a "Chain-of-Thought" (CoT) prompting strategy to encourage the model to reason through the data before arriving at a conclusion.

\subsubsection{Persona Adoption}
We explicitly define the model's role:
\begin{quote}
    "You are a conservative, risk-averse financial analyst with 20 years of experience. Your priority is capital preservation. You prefer to miss a profit opportunity rather than incur a loss."
\end{quote}
This persona instruction biases the model towards "HOLD" or "SELL" recommendations in uncertain markets.

\subsubsection{Context Injection}
We structure the input data as a JSON object within the prompt. This structured format reduces token usage compared to natural language descriptions and allows the model to parse the time-series data effectively.
\begin{lstlisting}[language=json, caption=Data Context Structure]
{
  "symbol": "AAPL",
  "current_price": 150.25,
  "history": [
    {"date": "2023-10-01", "close": 148.00, "volume": 50000000},
    ...
  ]
}
\end{lstlisting}

\subsubsection{Output Constraints}
To ensure the output is parsable by our application, we enforce a strict JSON schema using the OpenAI API's \texttt{response\_format} parameter.
\begin{lstlisting}[language=json, caption=Target Output Schema]
{
  "recommendation": "BUY" | "SELL" | "HOLD",
  "risk_level": "LOW" | "MEDIUM" | "HIGH",
  "reasoning": "string"
}
\end{lstlisting}

\subsection{Decision Logic}
The final trading signal is a composite of the quantitative filter and the qualitative LLM output.
\begin{itemize}
    \item \textbf{Step 1}: If $\sigma > Threshold$, return "High Volatility Warning" (No Trade).
    \item \textbf{Step 2}: Query LLM.
    \item \textbf{Step 3}: If LLM Recommendation is "BUY" AND Risk Level is "LOW", display "BUY".
    \item \textbf{Step 4}: Otherwise, display "HOLD" or "SELL".
\end{itemize}
This "double-check" mechanism ensures that the AI's hallucinations do not lead to reckless trading decisions.
